%!TEX root = ../main.tex
%----------------------------------------------------------------------
%  Section 7: Conclusion
%  Paper:   Theory-Informed Kernel Ridge Regression for Asset Pricing
%  Authors: Amedeo Andriollo and Juan F. Imbet
%----------------------------------------------------------------------

\section{Conclusion}\label{sec:conclusion}

Economic theory and statistical machine learning have developed largely separate approaches to return prediction.
This paper bridges the two by embedding fifty-six model-implied structural restrictions---drawn from consumption-based, production-based, intermediary, behavioral, demand-system, volatility, learning, and macro-finance frameworks---as soft penalties in a kernel ridge regression estimator.
Each restriction carries its own penalty multiplier $\mu_j$, selected by cross-validation, so the data determine which theories receive weight and which are shut down.
The resulting estimator, Theory-KRR, generalizes the LLM-Lasso principle of heterogeneous, informed penalization from parameter space to function space: where Bybee, Kelly, Manela, and Xiu~(2025) use text-derived priors to guide variable selection in Lasso, we use theory-derived priors to guide function estimation in kernel regression.

Theory-KRR outperforms standard KRR in out-of-sample return prediction, with $R^2_{\text{OOS}}$ gains of [XX.XX] percentage points and Sharpe ratio improvements of [XX.XX].
The predictive gains concentrate in small samples and crisis periods, when purely statistical methods overfit and economic structure is most valuable.
Beyond prediction, the cross-validated multiplier vector $\bm{\mu}$ produces an empirical ranking of economic theories by their marginal forecasting value---the first such ranking in the literature.
Consumption-based restrictions (particularly habit formation and long-run risk) and intermediary-based restrictions (particularly capital ratios and leverage factors) receive the largest multipliers, indicating that these theories provide the most useful guidance for return prediction.
Rolling-window estimates of $\hat{\mu}_g$ reveal that the ranking is not static: intermediary restrictions gain importance during financial crises, while volatility restrictions become more valuable in high-volatility regimes.

These findings carry implications along three dimensions.
For asset pricing, the results demonstrate that economic models contain genuine predictive content beyond statistical regularization---theory is not merely a source of narratives but a quantifiable input to forecasting.
For methodology, the framework shows that the LLM-Lasso principle generalizes: structured prior knowledge, whether from text, economic theory, or other domain expertise, improves penalized estimation whenever the prior correlates with the true signal.
For model evaluation, the cross-validated multipliers $\bm{\mu}$ provide a new lens for comparing asset-pricing theories---not by their in-sample fit to moments or their ability to rationalize anomalies, but by their marginal out-of-sample forecasting value conditional on all competing theories.

Several limitations define avenues for future work.
The structural parameters $\hat{\theta}_j$ entering each penalty are pre-estimated by GMM, introducing estimation error that propagates into the multiplier estimates; joint estimation of $(\hat{f}, \bm{\mu}, \bm{\theta})$ is a natural extension, though computationally demanding.
Our eight-group cross-validation aggregates heterogeneity within groups---estimating individual $\mu_j$ for all fifty-six restrictions with hierarchical priors would sharpen the ranking at the cost of a larger hyperparameter search.
The framework applies to kernel methods, but the same logic extends to neural networks: theory-informed regularization terms in the loss function of a deep network would combine the flexibility of neural architectures with the discipline of economic structure.
Finally, our empirical analysis covers U.S.\ equities; international markets, corporate bonds, and other asset classes present different challenges and different theoretical priors, and the relative ranking of theories may differ across these domains.
